% ============================================================
%  第五章（选作）：用 DSH 实现个人 Alpha —— 项目学习报告助手
%  （dsh-learning-report）
% ============================================================
\clearpage
\section{第五章（选作）：用 DSH 实现个人 Alpha —— 项目学习报告助手}

\subsection{背景：问题定义}

\textbf{一句话问题定义}：学习或接手一个陌生项目时，需要回答``它是什么、
怎么搭、怎么跑''——依赖清单、文档入口、目录结构、git 状态、构建命令，
人工要逐文件探测、逐个命令试跑，费时且容易漏。本 Alpha 把它做成一个
DSH 插件（\nolinkurl{code/ch5-alpha/dsh-plugins/} 仓库中的
\texttt{dsh-learning-report}）：用户一句话（``分析这个项目，写技术报告
和复现方案''），模型按插件分发的 skill 工作流自动完成——一次收集结构化
项目快照，产出 \texttt{TECHNICAL\_REPORT.md}（技术报告）与
\texttt{REPRODUCTION.md}（复现方案），全部落盘为文件。

插件已安装进\textbf{Windows 11 本机 DSH 的 desktop profile}
（\texttt{DSH\_HOME=C:\textbackslash Users\textbackslash KNtnx\textbackslash .dsh}）。

\subsection{实验目标}

\begin{enumerate}[leftmargin=4em,itemsep=0pt,topsep=2pt]
    \item 以 bundle 形式把插件装入本机 desktop profile，并验证
          \textbf{真实启动加载}（终端打印插件已加载）；
    \item 梳理插件的工具/skill 结构与数据来源（文件系统与 git），
          画出最小架构；
    \item 落实安全边界：只读收集、输出目录限定、敏感文件不读内容、
          人工确认；
    \item 记录一次真实输入输出（运行证据，见 5.4）。
\end{enumerate}

\subsection{步骤与环境}

环境：Windows 11（DSH\_HOME \texttt{C:\textbackslash Users\textbackslash KNtnx\textbackslash .dsh}，
desktop profile 为 web-app 形态）；DSH 0.1.0-rc.6（npm 全局）；插件
$@$0.1.0（本地 bundle）。插件仓库
\nolinkurl{code/ch5-alpha/dsh-plugins/}（pnpm workspaces；本 Alpha
采用项目报告插件 \texttt{dsh-learning-report}，其余包不在本 Alpha
范围）。安装方式与第 4 章相同（官方 bundle 流程：
\texttt{dsh plugin --profile desktop add <包 tgz>}），插件包内含
\texttt{src/learning-report.ts}（\texttt{collect\_project\_context}）、
\texttt{skills/learning-report/}（项目报告五阶段工作流）与
\texttt{lib/}、\texttt{cordis.patch.yml}。\textbf{真实加载证据}见图~\ref{fig:ch5_loaded}
（首次以默认端口启动时 \texttt{EADDRINUSE: 3080}，按提示改用
\texttt{--port 3081} 后加载成功——排错过程一并保留展示）：

\begin{figure}[H]
\centering
\includegraphics[width=0.95\textwidth]{ch5_plugin_loaded.jpg}
\caption{desktop profile 真实启动：\texttt{dsh --profile desktop
--port 3081} 后终端打印 \texttt{[project-learning-report] 插件已加载！}
（上方为端口占用排错记录；图中另有一行插件加载输出已按本章范围遮盖）。}
\label{fig:ch5_loaded}
\end{figure}

\subsection{结果：真实运行证据}

对真实项目 \textbf{HighWayFlow}（强化学习高速赛车 Gymnasium 环境项目，
位于 \texttt{E:\textbackslash game\textbackslash learn-llm}）执行一次
``写技术报告和复现方案''（图~\ref{fig:ch5_session}）：会话中加载
\texttt{project-learning-report} skill，
\texttt{collect\_project\_context} 首次调用完成结构化快照收集，
随后模型按工作流补读关键文件并落盘两份报告——会话统计共
\textbf{18 次工具调用}（1 次 \texttt{collect\_project\_context} +
补读/写盘等内置工具调用），随后汇报交付文件：

\begin{figure}[H]
\centering
\includegraphics[width=0.95\textwidth]{ch5_session_result.jpg}
\caption{运行结果（desktop 会话）：\texttt{project-learning-report
分析 HighWayFlow 项目}——会话共 18 次工具调用，输出目录
\texttt{E:\textbackslash game\textbackslash learn-llm\textbackslash
learning-highwayflow}，交付 \texttt{TECHNICAL\_REPORT.md}（328 行）与
\texttt{REPRODUCTION.md}（156 行），并给出技术报告与复现方案要点摘要。}
\label{fig:ch5_session}
\end{figure}

\subsubsection{交付文件 1：TECHNICAL\_REPORT.md}

按 SKILL 固定结构成文（项目概述/技术栈与依赖/目录结构/核心代码分析/
Git 状态/构建与运行），图~\ref{fig:ch5_tech} 为其中依赖与目录结构段
（依赖清单精确到版本与用途，目录树带中文批注——例如
\texttt{highwayflow/envs/vehicle/} 下 \texttt{kinematics.py}（车辆运动学）
与 \texttt{controller.py}（纵向/横向 PID 控制器））：

\begin{figure}[H]
\centering
\includegraphics[width=0.95\textwidth]{ch5_report_technical.jpg}
\caption{TECHNICAL\_REPORT.md（编辑器打开）：第 3 段 依赖与工具（gymnasium
0.31.0、numpy、pandas 等；testing/style 工具表；uv 作为包管理器）与
第 4 段 目录结构（带中文模块批注）——高亮行为截图中的表格例。}
\label{fig:ch5_tech}
\end{figure}

\subsubsection{交付文件 2：REPRODUCTION.md}

按 SKILL 固定结构成文（环境要求/依赖安装/复现步骤/预期结果/验证清单），
图~\ref{fig:ch5_repro} 为第 4 段复现结果与第 5 段验证清单（42 个测试 12.3s 全部
通过；验证清单逐项给出验证方法——如
\texttt{python -c "import highway\_env"} 预期输出 \texttt{1.22.1}——
与预期结果列一一对应）：

\begin{figure}[H]
\centering
\includegraphics[width=0.95\textwidth]{ch5_report_reproduction.jpg}
\caption{REPRODUCTION.md（编辑器打开）：第 4 段 复现结果（42 tests passed
in 12.3s、环境可视化、DQN 训练曲线）与第 5 段 验证清单（检查项/验证方法/
预期结果三列）。}
\label{fig:ch5_repro}
\end{figure}

至此完成``一句话 → 结构化快照 → 技术报告 + 复现方案''的完整闭环：
\texttt{collect\_project\_context} 一次调用拿到
manifest/README/AGENTS.md/文件清单/git 状态（省 token），模型仅补读
少数关键文件；两份报告全部落盘为文件（不在对话里贴长文）。

\subsection{分析与反思}

\textbf{最小架构}见图~\ref{fig:ch5_arch}（按已安装 bundle 的
\texttt{cordis.patch.yml}、工具声明与 SKILL.md 绘制）：用户一句话 →
skill 工作流阶段化驱动（阶段 1--5）→ 插件工具
\texttt{collect\_project\_context}（一次拿到结构化快照，省 token）与
宿主内置工具（read/write/grep/bash 补读与落盘）→ 数据/产物两路：
项目仓库（清单/文档/源码/git）$\Rightarrow$ 产物
\texttt{learning-<项目>/} 下两个 Markdown。skill 层负责``怎么编排''，
工具层负责``能做什么''，分离使工具可被其它工作流复用。

\begin{figure}[H]
\centering
\includegraphics[width=0.95\textwidth]{ch5_arch.png}
\caption{个人 Alpha 最小架构：用户/会话、skill 工作流、插件工具与
数据/产物四层；\texttt{collect\_project\_context} 与
\texttt{project-learning-report} 均来自已安装的 bundle。}
\label{fig:ch5_arch}
\end{figure}

\textbf{为什么适合用 DSH/插件化实现}：``一次收集 + 组织 + 落盘''是
\textbf{多工具、可组合、带权限边界}的职责——DSH 插件模型统一了参数校验
（\texttt{defineTool}）、输出 schema（供程序消费）与权限边界；skill
把收集-写作流程固化为不可变的会话行为，避免每次重复交代；bundle 可被
任意 profile 复用（已验证 desktop 形态加载）。

\textbf{安全边界}：
\begin{enumerate}[leftmargin=4em,itemsep=0pt,topsep=2pt]
    \item \textbf{只读收集}：\texttt{collect\_project\_context} 只读
          项目清单/文档/文件列表/git 状态（git 命令为只读形态），
          写操作限定在产物目录 \texttt{learning-<项目>/}；
    \item \textbf{敏感文件}：不读取 \texttt{.env} 等凭据文件（快照只含
          清单与文件名，不含文件内容密钥）；产物不含 API key；
    \item \textbf{不覆盖}：输出目录已存在时不覆盖已有报告，文件冲突
          采用追加或另命名（SKILL 规定）；\texttt{collect\_project\_context}
          返回的清单用于快照而非逐份粘贴全文（省 token 且保证
          上下文可控）；
    \item \textbf{人工确认}：报告/复现方案为模型生成的一手产出，
          交付前提示用户核对关键命令与预期结果。
\end{enumerate}

\textbf{局限性}与下一步：
\begin{enumerate}[leftmargin=4em,itemsep=0pt,topsep=2pt]
    \item 快照质量依赖源项目结构：manifest 识别覆盖 package.json、
          pyproject.toml、Cargo.toml、go.mod、requirements.txt，
          老式或混合项目可能识别不全；
    \item \texttt{git} 依赖本机 CLI（与第 4 章相同的环境前置）；
    \item 下一步：接入 DSH \texttt{tools/pre-execute} 做产出前检查；
          支持批量/定时学习任务；报告模板参数化。
\end{enumerate}

\begin{thebibliography}{9}

\bibitem{dsh-plugins-ch5}
本项目. dsh-plugins——个人插件仓库（项目学习报告插件）[EB/OL].
\nolinkurl{code/ch5-alpha/dsh-plugins/}，访问日期 2026-08-31.

\bibitem{dsh-skill-ch5}
DeepSeek. DSH 技能（skills）与插件开发文档[EB/OL].
\url{https://github.com/deepseek-ai/DeepSeek-Harness/tree/master/docs}，
访问日期 2026-08-31.

\end{thebibliography}
